% !TeX root = ../../Book3.tex
\section{Green 函数的构造}
\label{b3:sec:1703}

先在规则有界域构造核，再逐步放大区域。单调性与 Harnack 原理控制极限，最小性使构造不依赖所选穷竭。

\subsection{穷竭}
\label{b3:subsec:170301}

任意平面区域可取 \(D_j\) 光滑有界连通子域，\(\overline{D_j}\subset D_{j+1}\)，且并为 \(D\)。例如先取到补集距离至少 \(1/j\)、模不超过 \(j\) 的紧部分，添加连接路径，再在内部用有限圆盘和光滑边界逼近。极点包含在首项之后。

若 \(D_j\subset D_{j+1}\)，核之差在极点可去，且在 \(\partial D_j\) 非负；最大值原理给 \(g_{D_j}\le g_{D_{j+1}}\)。这是真正的单调量。不能直接说边界值逐项收敛，因为边界本身在移动。

\subsection{Perron 方法}
\label{b3:subsec:170302}

对于光滑有界 \(D_j\)，用 Perron 法解边值 \(\log|z-w|\)，得调和函数 \(h_j\)，令 \(g_j=-\log|z-w|+h_j\)。小圆上它趋正无穷、外边界为零，最大值原理给严格正性。

\begin{theorem}[Green 核的穷竭构造]\label{b3:thm:17-green-exhaustion}
若 \(g_j(z_0,w)\) 对某个 \(z_0\ne w\) 有界，则 \(g_j\) 局部一致收敛于有限正 Green 核；该核是所有具有同一对数奇性的正调和核中的最小者，因而与穷竭无关。
\end{theorem}
\begin{proof}
极点外用 Harnack 收敛得到调和极限。在极点周围固定一小圆，边界上的极限有界；\(h_j=g_j+\log|z-w|\) 在圆盘内调和，最大值原理把边界界传到内部。Harnack 或 Poisson 表示使 \(h_j\) 在较小圆盘收敛到调和函数，所以奇性保持。
任取竞争正核 \(G\)，\(G-g_j\) 在极点可去，且在 \(\partial D_j\) 非负，因此 \(g_j\le G\)。取极限得最小性。不同穷竭产生的有限核彼此作为竞争核，故相等。若一个穷竭发散而另一个有限，后者又给前者共同上界，矛盾。
\end{proof}
有限核的对称性由光滑有界核的对称性取极限得到。

\subsection{圆盘与半平面}
\label{b3:subsec:170303}

圆盘与上半平面的显式核为
\[
g_{B(0,R)}(z,w)=\log\left|\frac{R^2-\overline wz}{R(z-w)}\right|,
\quad
g_{\mathfrak H}(z,w)=\log\left|\frac{z-\overline w}{z-w}\right|.
\]
分子零点在区域外，故余项调和。边界上分子和分母的模相等，且区域内两模之比大于一。圆盘由唯一性，半平面由共形变换和最小性确认。

极点在圆盘中心时核简化成 \(\log(R/|z|)\)。上半平面公式是负对数点源与反射源之差；反射点必须为 \(\overline w\)，不能给反射源保留同号。

\subsection{比较原理}
\label{b3:subsec:170304}

\begin{proposition}\label{b3:prop:17-domain-monotonicity}
若 Green 域 \(D\subset E\)，则 \(g_D(z,w)\le g_E(z,w)\)。
\end{proposition}
\begin{proof}
\(E\) 的核限制到 \(D\) 是一个正竞争核，具有正确奇性，最小性立即给不等式。
\end{proof}
将边界推远会减弱零势约束，核因而增大。对有界域，这个原则用包围圆盘给有限上界；它也使某些无界域可由半平面或其他已知 Green 域控制。


\ProseParagraphBreak

穷竭极限还保留极点参数的局部稳定性。在两个不相交的相对紧小圆盘 \(U,V\Subset D\) 上，光滑穷竭核在两个变量各自调和且对称。固定源点先用 Harnack 控制观察点，再利用对称性在源点方向使用同一估计，就可在 \(U\times V\) 上得到共同上界及导数界。极限因此在远离对角线处连续，并在每个变量中调和。仅仅知道逐点单调极限有限，还不能未经此控制就交换源点极限或把核当成连续的双变量函数。

\ProseParagraphBreak

在圆盘里，核的正性可由一个恒等式直接核对：
\[
|R^2-\overline wz|^2-R^2|z-w|^2=(R^2-|z|^2)(R^2-|w|^2)>0.
\]
两点均在圆盘内部时右端严格正，故模之比大于一。上半平面的相应恒等式是 \(|z-\overline w|^2-|z-w|^2=4\operatorname{Im}z\operatorname{Im}w>0\)。这些实代数算式避免凭图形猜测反射点的位置和源的符号；也说明两核在光滑边界分别降到零。

\ProseParagraphBreak

若 \(D\subsetneq E\) 且两域连通，区域单调性一般是严格的。在 \(D\) 内两核差调和、非负；如果在一个内部非极点位置为零，强最小值原理使差在 \(D\) 恒零。对光滑域，如果 \(\partial D\) 有一个点位于 \(E\) 内，\(E\) 的核在那里严格正而 \(D\) 的核取零，便产生矛盾。对于粗糙域，删去一个极性孔可能不改变核，所以不能仅由集合严格包含就断言严格不等式。严格性需要边界差异在势论中具有实际影响。
